% GENERATED FILE. Edit canonical lessons, then run npm run generate:books.
% canonical-source-hash: fnv1a64:c81ec8861928b933
% canonical-lessons: BN-C524-ijel, BN-C524-jolrong, BN-C524-noksha, BN-C524-khodai-kora, BN-C524-rangano

\chapter{Paint and Carving}
\label{ch:bn-paint-and-carving}
\hlchaptermodality{524}

\begin{chapteropening}
\textbf{By the end of this chapter:} \emph{I can talk about brushes, watercolour, patterns and carving.}

Complete the last lesson of chapter 524: I can talk about brushes, watercolour, patterns and carving.
\end{chapteropening}

\section[\texorpdfstring{\bn{ইজেল} (ijel)}{ইজেল (ijel)}]{\bn{ইজেল} (ijel) — an easel}

\label{lesson:BN-C524-ijel}



\begin{quote}
Before the new one: say the Bengali for a fish-hook, then the Bengali for to swim across.
\end{quote}

\subsection*{You'll want to know: \bn{ইজেল}}
\textbf{\bn{ইজেল}} — \emph{ijel} — \textquotedblleft{}an easel\textquotedblright{}.

\begin{quote}
\bn{শিল্পী} \bn{ইজেলে} \bn{একটা} \bn{নতুন} \bn{ছবি} \bn{রাখলেন।}
\end{quote}

\begin{quote}
\emph{Shilpī ijele ekṭā nôtun chhôbi rākhlen.}
\end{quote}

\begin{quote}
The artist put a new picture on the easel.
\end{quote}

\begin{cousinweb}
\textbf{\bn{ইজেল}} (\emph{ijel}) is borrowed from English \emph{easel}.
\end{cousinweb}

\begin{grammarlens}[title={What you've built}]
One more word for describing what you lived and felt.
\end{grammarlens}

\subsection*{Your turn}
\begin{itemize}
\raggedright
  \item \emph{Say it:} \emph{ijel}
  \item \emph{Say it:} \emph{ijel}, once more
  \item \emph{Say it:} \emph{Shilpī ijele ekṭā nôtun chhôbi rākhlen.} — The artist put a new picture on the easel.
  \item \emph{Recall:} say the Bengali for a fish-hook, then the Bengali for to swim across, then say \emph{ijel} again
\end{itemize}

\begin{tcolorbox}[breakable,colback=teal!4,colframe=teal!35!black,title={Before you move on}]
What does \bn{ইজেল} mean? (\textquotedblleft{}An easel\textquotedblright{}.) Say it once more.
\end{tcolorbox}
\section[\texorpdfstring{\bn{জলরং} (jôlrông)}{জলরং (jôlrông)}]{\bn{জলরং} (jôlrông) — watercolour}

\label{lesson:BN-C524-jolrong}



\begin{quote}
Before the new one: say the Bengali for to swim across, then the Bengali for an easel.
\end{quote}

\subsection*{You'll want to know: \bn{জলরং}}
\textbf{\bn{জলরং}} — \emph{jôlrông} — \textquotedblleft{}watercolour\textquotedblright{}.

\begin{quote}
\bn{শিল্পী} \bn{জলরঙে} \bn{একটা} \bn{নদীর} \bn{ছবি} \bn{আঁকলেন।}
\end{quote}

\begin{quote}
\emph{Shilpī jôlrônge ekṭā nôdīr chhôbi ā̃klen.}
\end{quote}

\begin{quote}
The artist painted a river in watercolour.
\end{quote}

Before a vowel ending \bn{ং} is written \bn{ঙ}: \bn{জলরঙে}, in watercolour.

\begin{cousinweb}
\textbf{\bn{জলরং}} (\emph{jôlrông}) joins \bn{জল}, \textquotedblleft{}water\textquotedblright{}, to \bn{রং}, \textquotedblleft{}colour\textquotedblright{}, a word that matches both Sanskrit \emph{raṅga} (CDIAL 10560) and Persian \emph{rang}: colour mixed with water.
\end{cousinweb}

\begin{grammarlens}[title={What you've built}]
One more word for describing what you lived and felt.
\end{grammarlens}

\subsection*{Your turn}
\begin{itemize}
\raggedright
  \item \emph{Say it:} \emph{jôlrông}
  \item \emph{Say it:} \emph{jôlrông}, once more
  \item \emph{Say it:} \emph{Shilpī jôlrônge ekṭā nôdīr chhôbi ā̃klen.} — The artist painted a river in watercolour.
  \item \emph{Recall:} say the Bengali for to swim across, then the Bengali for an easel, then say \emph{jôlrông} again
\end{itemize}

\begin{tcolorbox}[breakable,colback=teal!4,colframe=teal!35!black,title={Before you move on}]
What does \bn{জলরং} mean? (\textquotedblleft{}Watercolour\textquotedblright{}.) Say it once more.
\end{tcolorbox}
\section[\texorpdfstring{\bn{নকশা} (nôkshā)}{নকশা (nôkshā)}]{\bn{নকশা} (nôkshā) — a design; a pattern}

\label{lesson:BN-C524-noksha}



\begin{quote}
Before the new one: say the Bengali for an easel, then the Bengali for watercolour.
\end{quote}

\subsection*{You'll want to know: \bn{নকশা}}
\textbf{\bn{নকশা}} — \emph{nôkshā} — \textquotedblleft{}a design; a pattern\textquotedblright{}.

\begin{quote}
\bn{শাড়িটার} \bn{নকশা} \bn{সুন্দর।}
\end{quote}

\begin{quote}
\emph{Shāṛiṭār nôkshā sundôr.}
\end{quote}

\begin{quote}
The sari's pattern is beautiful.
\end{quote}

\begin{cousinweb}
\textbf{\bn{নকশা}} (\emph{nôkshā}) is usually traced to Arabic \emph{naqsh}, \textquotedblleft{}a drawing, an engraving\textquotedblright{}, through Persian \emph{naqsha}.
\end{cousinweb}

\begin{grammarlens}[title={What you've built}]
One more word for describing what you lived and felt.
\end{grammarlens}

\subsection*{Your turn}
\begin{itemize}
\raggedright
  \item \emph{Say it:} \emph{nôkshā}
  \item \emph{Say it:} \emph{nôkshā}, once more
  \item \emph{Say it:} \emph{Shāṛiṭār nôkshā sundôr.} — The sari's pattern is beautiful.
  \item \emph{Recall:} say the Bengali for an easel, then the Bengali for watercolour, then say \emph{nôkshā} again
\end{itemize}

\begin{tcolorbox}[breakable,colback=teal!4,colframe=teal!35!black,title={Before you move on}]
What does \bn{নকশা} mean? (\textquotedblleft{}A design; a pattern\textquotedblright{}.) Say it once more.
\end{tcolorbox}
\section[\texorpdfstring{\bn{খোদাই} \bn{করা} (khodāi kôrā)}{খোদাই করা (khodāi kôrā)}]{\bn{খোদাই} \bn{করা} (khodāi kôrā) — to carve; to engrave}

\label{lesson:BN-C524-khodai-kora}



\begin{quote}
Before the new one: say the Bengali for watercolour, then the Bengali for a pattern.
\end{quote}

\subsection*{You'll want to know: \bn{খোদাই} \bn{করা}}
\textbf{\bn{খোদাই} \bn{করা}} — \emph{khodāi kôrā} — \textquotedblleft{}to carve; to engrave\textquotedblright{}.

\begin{quote}
\bn{কারিগর} \bn{কাঠে} \bn{একটা} \bn{ফুল} \bn{খোদাই} \bn{করল।}
\end{quote}

\begin{quote}
\emph{Kārigôr kāṭhe ekṭā phul khodāi kôrlô.}
\end{quote}

\begin{quote}
The craftsman carved a flower into the wood.
\end{quote}

Past (he/she …): \bn{তিনি} \bn{পাথরে} \bn{নাম} \bn{খোদাই} \bn{করলেন}, \emph{tini pāthôre nām khodāi kôrlen}.

\begin{cousinweb}
\textbf{\bn{খোদাই} \bn{করা}} (\emph{khodāi kôrā}) joins \bn{খোদাই}, \textquotedblleft{}carving\textquotedblright{}, built on a verb CDIAL lists as Bengali \emph{khodā}, \textquotedblleft{}to dig, to carve\textquotedblright{}, from a reconstructed \emph{khōdd}, \textquotedblleft{}dig\textquotedblright{} (CDIAL 3934), to \bn{করা}, \textquotedblleft{}to do\textquotedblright{}.
\end{cousinweb}

\begin{grammarlens}[title={What you've built}]
One more word for describing what you lived and felt.
\end{grammarlens}

\subsection*{Your turn}
\begin{itemize}
\raggedright
  \item \emph{Say it:} \emph{khodāi kôrā}
  \item \emph{Say it:} \emph{khodāi kôrā}, once more
  \item \emph{Say it:} \emph{Kārigôr kāṭhe ekṭā phul khodāi kôrlô.} — The craftsman carved a flower into the wood.
  \item \emph{Recall:} say the Bengali for watercolour, then the Bengali for a pattern, then say \emph{khodāi kôrā} again
\end{itemize}

\begin{tcolorbox}[breakable,colback=teal!4,colframe=teal!35!black,title={Before you move on}]
What does \bn{খোদাই} \bn{করা} mean? (\textquotedblleft{}To carve; to engrave\textquotedblright{}.) Say it once more.
\end{tcolorbox}
\section[\texorpdfstring{\bn{রাঙানো} (rāngāno)}{রাঙানো (rāngāno)}]{\bn{রাঙানো} (rāngāno) — to colour; to dye}

\label{lesson:BN-C524-rangano}



\begin{quote}
Before the new one: say the Bengali for a pattern, then the Bengali for to carve.
\end{quote}

\subsection*{You'll want to know: \bn{রাঙানো}}
\textbf{\bn{রাঙানো}} — \emph{rāngāno} — \textquotedblleft{}to colour; to dye\textquotedblright{}.

\begin{quote}
\bn{শিল্পী} \bn{নতুন} \bn{বাড়ির} \bn{দেওয়াল} \bn{রাঙালেন।}
\end{quote}

\begin{quote}
\emph{Shilpī nôtun bāṛir deoyāl rāngālen.}
\end{quote}

\begin{quote}
The artist coloured the walls of the new house.
\end{quote}

Past (he/she …): \bn{তিনি} \bn{কাপড়টা} \bn{রাঙালেন}, \emph{tini kāpṛṭā rāngālen}.

\begin{cousinweb}
\textbf{\bn{রাঙানো}} (\emph{rāngāno}) comes from a reconstructed \emph{raṅgayati}, \textquotedblleft{}dyes\textquotedblright{}, built on Sanskrit \emph{raṅga}, \textquotedblleft{}colour\textquotedblright{}; CDIAL lists Bengali \emph{rāṅgāna}, \textquotedblleft{}to redden\textquotedblright{} (CDIAL 10570).
\end{cousinweb}

\begin{grammarlens}[title={What you've built}]
One more word for describing what you lived and felt.
\end{grammarlens}

\subsection*{Your turn}
\begin{itemize}
\raggedright
  \item \emph{Say it:} \emph{rāngāno}
  \item \emph{Say it:} \emph{rāngāno}, once more
  \item \emph{Say it:} \emph{Shilpī nôtun bāṛir deoyāl rāngālen.} — The artist coloured the walls of the new house.
  \item \emph{Recall:} say the Bengali for a pattern, then the Bengali for to carve, then say \emph{rāngāno} again
\end{itemize}

\begin{tcolorbox}[breakable,colback=teal!4,colframe=teal!35!black,title={Before you move on}]
What does \bn{রাঙানো} mean? (\textquotedblleft{}To colour; to dye\textquotedblright{}.) Say it once more.
\end{tcolorbox}
